\section{System Architecture and Modeling}\label{sec:problem_system_model}
\noindent The Lightweight Blockchain-Based Data Trading (LBDT) ecosystem operates within a multi-zone urban Internet of Vehicles (IoV) environment. The system models three primary entities:

\begin{enumerate}
    \item \textbf{Vehicles ($V = \{v_1, v_2, \dots, v_M\}$):} Mobile nodes equipped with On-Board Units (OBUs) and wireless communication interfaces (IEEE 802.11p DSRC / C-V2X). Vehicles act dynamically as \textbf{Data Buyers} (seeking real-time traffic updates, sensor feeds, or road hazard alerts) or \textbf{Data Sellers} (offering verified sensor payloads for financial rewards).
    \item \textbf{Roadside Units ($R = \{r_1, r_2, \dots, r_K\}$):} Fixed roadside infrastructure nodes located along intersections and highway segments. RSUs possess higher computational, storage, and networking capacities. RSUs act as \textbf{Zone Authority Nodes}, hosting consensus committees, executing double auctions, and maintaining local ledger states.
    \item \textbf{Geographic Zones ($Z = \{z_1, z_2, \dots, z_P\}$):} Dynamically sharded spatial partitions covering the urban road network. Each zone $z_p$ is managed by a subset of localized RSUs and participating vehicular consensus nodes.
\end{enumerate}

\section{Security Threat Model and Adversary Assumptions}\label{sec:problem_threat_model}
\noindent In an open and decentralized IoV environment, we consider a hostile network environment where up to $f < N/3$ consensus nodes (and arbitrary vehicular participants) may be compromised or malicious. We classify potential attack vectors as follows:

\begin{itemize}
    \item \textbf{Data Tampering and Fraudulent Sellers:} Malicious sellers transmit corrupted, outdated, or fabricated sensor payloads after winning auction bids to harvest profits without incurring data acquisition costs.
    \item \textbf{Dishonest Bidding and Free-Riding Buyers:} Malicious buyers submit artificial bid prices ($BP$) or attempt to reuse single data purchases across unauthorized third-party vehicles without payment.
    \item \textbf{Selective Packet Dropping and Sybil Attacks:} Adversarial nodes deliberately drop relay packets, suppress consensus messages, or spawn multiple pseudonymous identities to distort reputation scores.
    \item \textbf{BFT Committee Corruption and Leader Bias:} Colluding RSUs attempt to bias committee leader selection, delay transaction confirmation, or insert invalid blocks into the ledger.
\end{itemize}

\section{System Objectives and Optimization Goals}\label{sec:problem_objectives}
\noindent The LBDT scheme has developed five major optimization goals to respond to these security problems and operational constraints, namely:
\begin{enumerate}
    \item \textbf{Latency Optimization:} Ensure less than 3 seconds end-to-end transaction confirmation latency for high-frequency trading operations in conditions of high vehicle density.
    \item \textbf{Reputation Accuracy:} Track the past interactions between nodes and apply penalties for dishonest behavior as per old Gompertz decay principle supported by machine learning ($p_{honest}$) so that malicious nodes are eliminated before entering auction matching.
    \item \textbf{Economic Fairness and Truthfulness:} Ensure that the price of the double auction clearing ($p_{win}$) reflects data freshness and seller reliability by implementing bid loss penalties.
    \item \textbf{Consensus Efficiency and Safety:}Ensure linear communication overhead through BFT of (5N-3) messages per round together with committee safety guarantee (malicious reputation ratio less than $< 33.3\%$).
    \item \textbf{Dynamic Load Balancing:} Ensure optimal utilization of zones ($U_z$) by  sharding split and merge operations.
\end{enumerate}
